            因此 Jordan 形式为  
            \begin{equation*}
              \begin{mat}[r]
                -2  &0  &0  \\
                 0  &4  &0  \\
                 0  &1  &4
              \end{mat}
            \end{equation*}
            而一个合适的基如下。
            \begin{equation*}
              B=\cat{B_{-2}}{B_4}=
                \sequence{\colvec[r]{1 \\ 1 \\ 1},
                          \colvec[r]{0 \\ -1 \\ 1},
                          \colvec[r]{-1 \\ 1 \\ -1}}
            \end{equation*}
        \partsitem 该矩阵的特征多项式
            为 \( c(x)=(2-x)^3=-1\cdot (x-2)^3 \)。
            这个矩阵只有一个特征值 $\lambda=2$。
            通过求 $T-2I$ 的各次幂可得
            \begin{equation*}
              \nullspace{t-2}=\set{\colvec{-y \\ y \\ 0}
                                      \suchthat y\in\C} 
              \qquad
              \nullspace{(t-2)^2}=\set{\colvec{-y-(1/2)z \\ y \\ z}
                                          \suchthat y,z\in\C}
            \end{equation*}
            以及
            \begin{equation*} 
              \nullspace{(t-2)^3}=\C^3
            \end{equation*}
            所以
            $t-2$ 在一个相伴的
            链基上的作用是
            $\vec{\beta}_1\mapsto\vec{\beta}_2\mapsto
                   \vec{\beta}_3\mapsto\zero$。
            Jordan 形式为
            \begin{equation*}
                  \begin{mat}[r]
                    2  &0  &0  \\
                    1  &2  &0  \\
                    0  &1  &2
                  \end{mat}
            \end{equation*}
            而一个可选的基如下。
            \begin{equation*}
              B=\sequence{\colvec[r]{0 \\ 1 \\ 0},
                          \colvec[r]{7 \\ -9 \\ 4},
                          \colvec[r]{-2 \\ 2 \\ 0}}
            \end{equation*}
        \partsitem 特征多项式
            \( c(x)=(1-x)^3=-(x-1)^3 \) 只有一个单根，
            因此该矩阵只有一个特征值 $\lambda=1$。
            求 $T-1I$ 的各次幂
            并计算零空间
            \begin{equation*}
               \nullspace{t-1}=\set{\colvec{-2y+z \\ y \\ z}
                                      \suchthat y,z\in\C} 
               \qquad
               \nullspace{(t-1)^2}=\C^3 
            \end{equation*}
            可知幂零映射 $t-1$ 在一个链基上
            的作用是
            $\vec{\beta}_1\mapsto\vec{\beta}_2\mapsto\zero$ 与
            $\vec{\beta}_3\mapsto\zero$。
            因此 Jordan 形式为
            \begin{equation*}
                  J=
                  \begin{mat}[r]
                    1  &0  &0  \\
                    1  &1  &0  \\
                    0  &0  &1
                  \end{mat}
            \end{equation*}
            而一个合适的基（与 $t-1$ 相伴的
            链基）如下。
            \begin{equation*}
              B=\sequence{\colvec[r]{0 \\ 1 \\ 0},
                          \colvec[r]{2 \\ -2 \\ -2},
                          \colvec[r]{1 \\ 0 \\ 1}}
            \end{equation*}
        \partsitem 特征多项式用手算
            有点大，但尚可应付：
            \( c(x)=x^4-24x^3+216x^2-864x+1296=(x-6)^4 \)。
            这是单特征值映射，所以
            变换 $t-6$ 是幂零的。
            零空间
            \begin{equation*}
               \nullspace{t-6}=\set{\colvec{-z-w \\ -z-w \\ z \\ w}
                                      \suchthat z,w\in\C} 
               \qquad
               \nullspace{(t-6)^2}=\set{\colvec{x \\ -z-w \\ z \\ w}
                                      \suchthat x,z,w\in\C} 
            \end{equation*}
            以及
            \begin{equation*}
               \nullspace{(t-6)^3}=\C^4 
            \end{equation*}
            连同各零度一起表明
            $t-6$ 在一个链基上的作用是
            $\vec{\beta}_1\mapsto\vec{\beta}_2\mapsto
                   \vec{\beta}_3\mapsto\zero$ 与
            $\vec{\beta}_4\mapsto\zero$。
            Jordan 形式为
            \begin{equation*}
              \begin{mat}[r]
                6  &0  &0  &0  \\
                1  &6  &0  &0  \\
                0  &1  &6  &0  \\
                0  &0  &0  &6  \\
              \end{mat}
            \end{equation*}
            而找一个合适的链基是例行的操作。
            \begin{equation*}
              B=\sequence{\colvec[r]{0 \\ 0 \\ 0 \\ 1},
                          \colvec[r]{2 \\ -1 \\ -1 \\ 2},
                          \colvec[r]{3 \\ 3 \\ -6 \\ 3},
                          \colvec[r]{-1 \\ -1 \\ 1 \\ 0}}
            \end{equation*}
      \end{exparts}  
    \end{answer}
  \recommended \item
    求特征多项式为 \( (x-1)^2(x+2)^2  \) 的
    变换的所有可能的 Jordan 形式。
    \begin{answer}
      有两个特征值，$\lambda_1=-2$ 与 $\lambda_2=1$。
      $t+2$ 限制在
      $\gennullspace{t+2}$ 上，在一个相伴的
      链基上可能有下面两种作用。
      \begin{equation*}
        \vec{\beta}_1\mapsto\vec{\beta}_2\mapsto\zero
        \qquad
        \begin{array}[t]{l}
          \vec{\beta}_1\mapsto\zero  \\
          \vec{\beta}_2\mapsto\zero 
        \end{array}
      \end{equation*}
      $t-1$ 限制在
      $\gennullspace{t-1}$ 上，在一个相伴的
      链基上可能有下面两种作用。
      \begin{equation*}
        \vec{\beta}_3\mapsto\vec{\beta}_4\mapsto\zero
        \qquad
        \begin{array}[t]{l}
          \vec{\beta}_3\mapsto\zero  \\
          \vec{\beta}_4\mapsto\zero 
        \end{array}
      \end{equation*}
      合起来共有四种可能的 Jordan 形式：
      两个都取第一种作用，第二种与第一种，第一种与第二种，
      以及两个都取第二种作用。
      \begin{equation*}
        \begin{mat}[r]
          -2  &0  &0  &0  \\
           1  &-2 &0  &0  \\
           0  &0  &1  &0  \\
           0  &0  &1  &1
        \end{mat}
        \quad
        \begin{mat}[r]
          -2  &0  &0  &0  \\
           0  &-2 &0  &0  \\
           0  &0  &1  &0  \\
           0  &0  &1  &1
        \end{mat}
        \quad
        \begin{mat}[r]
          -2  &0  &0  &0  \\
           1  &-2 &0  &0  \\
           0  &0  &1  &0  \\
           0  &0  &0  &1
        \end{mat}
        \quad
        \begin{mat}[r]
          -2  &0  &0  &0  \\
           0  &-2 &0  &0  \\
           0  &0  &1  &0  \\
           0  &0  &0  &1
        \end{mat}
     \end{equation*}  
    \end{answer}
  \item 
    求特征多项式为 \( (x-1)^3(x+2) \) 的
    变换的所有可能的 Jordan 形式。
    \begin{answer}
     $t+2$ 限制在
     $\gennullspace{t+2}$ 上只能有作用
     $\vec{\beta}_1\mapsto\zero$。
     $t-1$ 限制在 $\gennullspace{t-1}$ 上，在一个相伴的链基上可能有
     下面三种作用中的任意一种。
     \begin{equation*}
        \vec{\beta}_2\mapsto\vec{\beta}_3\mapsto\vec{\beta}_4\mapsto\zero
        \qquad
        \begin{array}[t]{l}
          \vec{\beta}_2\mapsto\vec{\beta}_3\mapsto\zero  \\
          \vec{\beta}_4\mapsto\zero 
        \end{array}
        \qquad
        \begin{array}[t]{l}
          \vec{\beta}_2\mapsto\zero  \\
          \vec{\beta}_3\mapsto\zero  \\
          \vec{\beta}_4\mapsto\zero 
        \end{array}
     \end{equation*}
     合起来共有三种可能的 Jordan 形式：
     由 $t-1$ 的第一种作用（连同 $t+2$ 唯一的
     作用）产生的形式、由第二种作用产生的形式，
     以及由第三种作用产生的形式。
     \begin{equation*}
       \begin{mat}[r]
         -2  &0  &0  &0  \\
          0  &1  &0  &0  \\
          0  &1  &1  &0  \\
          0  &0  &1  &1
       \end{mat}
       \quad
       \begin{mat}[r]
         -2  &0  &0  &0  \\
          0  &1  &0  &0  \\
          0  &1  &1  &0  \\
          0  &0  &0  &1
       \end{mat}
       \quad
       \begin{mat}[r]
         -2  &0  &0  &0  \\
          0  &1  &0  &0  \\
          0  &0  &1  &0  \\
          0  &0  &0  &1
       \end{mat}
     \end{equation*}
    \end{answer}
  \recommended \item
    求特征多项式为 \( (x-2)^3(x+1) \)、极小多项式为
    \( (x-2)^2(x+1) \) 的变换的所有可能的 Jordan 形式。
    \begin{answer}
      $t+1$ 在 $\gennullspace{t+1}$ 的一个链基上的作用
      必为 $\vec{\beta}_1\mapsto\zero$。 
      由于极小多项式中 \( x-2 \) 的幂次，$t-2$ 的一个链基
      长度为二，因此
      \( t-2 \) 在 \( \gennullspace{t-2} \) 上的作用
      必须具有如下形式。
      \begin{equation*}
        \begin{array}[t]{l}
          \vec{\beta}_2\mapsto\vec{\beta}_3\mapsto\zero  \\
          \vec{\beta}_4\mapsto\zero 
        \end{array}        
      \end{equation*}
      因此只可能有一种 Jordan 形式。
      \begin{equation*}
          \begin{mat}[r]
            -1  &0  &0  &0  \\
             0  &2  &0  &0  \\
             0  &1  &2  &0  \\
             0  &0  &0  &2
          \end{mat}
       \end{equation*}
      \end{answer}
  \item 
    求特征多项式为 \( (x-2)^4(x+1) \)、极小多项式为
    \( (x-2)^2(x+1) \) 的变换的所有可能的 Jordan 形式。
    \begin{answer}
      有两种可能的 Jordan 形式。
      $t+1$ 在 $\gennullspace{t+1}$ 的一个链基上的作用
      必为 $\vec{\beta}_1\mapsto\zero$。
      对于 $t-2$ 在
      $\gennullspace{t-2}$ 的一个链基上的作用，在这个特征
      多项式与极小多项式之下有两种可能的情形。
      \begin{equation*}
        \begin{array}[t]{l}
          \vec{\beta}_2\mapsto\vec{\beta}_3\mapsto\zero  \\
          \vec{\beta}_4\mapsto\vec{\beta}_5\mapsto\zero  
        \end{array}        
        \qquad
        \begin{array}[t]{l}
          \vec{\beta}_2\mapsto\vec{\beta}_3\mapsto\zero  \\
          \vec{\beta}_4\mapsto\zero                      \\
          \vec{\beta}_5\mapsto\zero                      
        \end{array}        
      \end{equation*}
      由此得到的 Jordan 形式矩阵如下。 
      \begin{equation*}
        \begin{mat}[r]
          -1  &0  &0  &0  &0  \\
           0  &2  &0  &0  &0  \\
           0  &1  &2  &0  &0  \\
           0  &0  &0  &2  &0  \\
           0  &0  &0  &1  &2
        \end{mat}
        \qquad
        \begin{mat}[r]
          -1  &0  &0  &0  &0  \\
           0  &2  &0  &0  &0  \\
           0  &1  &2  &0  &0  \\
           0  &0  &0  &2  &0  \\
           0  &0  &0  &0  &2
        \end{mat}
     \end{equation*}  
    \end{answer}
  \recommended \item 将它们对角化。
    \begin{exparts*}
       \partsitem \( \begin{mat}[r]
                  1  &1  \\
                  0  &0
                \end{mat}  \)
       \partsitem \( \begin{mat}[r]
                  0  &1  \\
                  1  &0
                \end{mat}  \)
    \end{exparts*}
    \begin{answer}
      \begin{exparts}
        \partsitem 特征多项式为 \( c(x)=x(x-1) \)。
          对于 $\lambda_1=0$，有
          \begin{equation*}
            \nullspace{t-0}=\set{\colvec{-y \\ y}
                                 \suchthat y\in\C }  
          \end{equation*} 
          （当然，$t^2$ 的零空间与之相同）。
          对于 $\lambda_2=1$，
          \begin{equation*}
            \nullspace{t-1}=\set{\colvec{x \\ 0}
                                     \suchthat x\in\C }  
          \end{equation*}
          （而 $(t-1)^2$ 的零空间与之相同）。
          我们可以取这个基
          \begin{equation*}
            B=\sequence{\colvec[r]{1 \\ -1},\colvec[r]{1 \\ 0}}
          \end{equation*}
          来得到对角化。
          \begin{equation*}
            \begin{mat}[r]
              1  &1  \\
             -1  &0
            \end{mat}^{-1}
            \begin{mat}[r]
              1  &1  \\
              0  &0
            \end{mat}
            \begin{mat}[r]
              1  &1  \\
             -1  &0
            \end{mat}
            =
            \begin{mat}[r]
              0  &0  \\
              0  &1
            \end{mat}
          \end{equation*}
        \partsitem 特征多项式为
          \( c(x)=x^2-1=(x+1)(x-1) \)。
          对于 $\lambda_1=-1$，
          \begin{equation*}
            \nullspace{t+1}=\set{\colvec{-y \\ y}
                                    \suchthat y\in\C } 
          \end{equation*}
          且 $(t+1)^2$ 的零空间与之相同。
          对于 $\lambda_2=1$，
          \begin{equation*}
            \nullspace{t-1}=\set{\colvec{y \\ y}
                                    \suchthat y\in\C } 
          \end{equation*}
          且 $(t-1)^2$ 的零空间与之相同。
          我们可以取这个基
          \begin{equation*}
            B=\sequence{\colvec[r]{1 \\ -1},\colvec[r]{1 \\ 1}}
          \end{equation*}
          来得到一个对角化。
          \begin{equation*}
            \begin{mat}[r]
              1  &1  \\
              1  &-1
            \end{mat}^{-1}
            \begin{mat}[r]
              0  &1  \\
              1  &0
            \end{mat}
            \begin{mat}[r]
              1   &1  \\
              -1  &1
            \end{mat}
            =
            \begin{mat}[r]
              -1  &0  \\
              0   &1
            \end{mat}
          \end{equation*}
      \end{exparts}  
     \end{answer}
  \recommended \item 
    求表示 \( \polyspace_3 \) 上微分算子的
    Jordan 矩阵。
    \begin{answer}
      变换 $\map{d/dx}{\polyspace_3}{\polyspace_3}$
      是幂零的。
      它在基 \( B=\sequence{x^3,3x^2,6x,6} \) 上的作用
      是 $x^3\mapsto 3x^2\mapsto 6x\mapsto 6\mapsto 0$。
      它的 Jordan 形式就是它作为幂零矩阵的标准形。
      \begin{equation*}
         J=
         \begin{mat}[r]
           0  &0  &0  &0  \\
           1  &0  &0  &0  \\
           0  &1  &0  &0  \\
           0  &0  &1  &0
         \end{mat}
      \end{equation*}
    \end{answer}
   \recommended \item 
      判断这两个矩阵是否相似。
      \begin{equation*}
         \begin{mat}[r]
            1  &-1 \\
            4  &-3 \\
         \end{mat}
         \qquad
         \begin{mat}[r]
           -1  &0  \\
            1  &-1 \\
         \end{mat}
      \end{equation*}
      \begin{answer}
        相似。
        两者都具有特征多项式 $(x+1)^2$。
        计算 $T_1+1\cdot I$ 与
        $T_2+1\cdot I$ 的各次幂得到下面两个。
        \begin{equation*}
          \nullspace{t_1+1}=\set{\colvec{y/2 \\ y} \suchthat y\in\C}
          \qquad
          \nullspace{t_2+1}=\set{\colvec{0 \\ y} \suchthat y\in\C}
        \end{equation*}
        （当然，对每一个来说，其平方的零空间
        都是整个空间。）
        零度的增长方式表明每一个都
        相似于这个 Jordan 形式矩阵
        \begin{equation*}
           \begin{mat}[r]
             -1  &0  \\
              1  &-1 \\
           \end{mat}
        \end{equation*}
        因而它们彼此相似。  
      \end{answer}
  \item 
     求这个矩阵的 Jordan 形式。
     \begin{equation*}
        \begin{mat}[r]
           0  &-1  \\
           1  &0
        \end{mat}
     \end{equation*}
     同时给出一个 Jordan 基。
     \begin{answer}
       它的特征多项式是
       \( c(x)=x^2+1 \)，它有复根
       \( x^2+1=(x+i)(x-i) \)。
       因为根互不相同，
       所以该矩阵可对角化，其 Jordan 形式就是那个
       对角矩阵。 
       \begin{equation*}
         \begin{mat}[r]
           -i  &0  \\
            0  &i
         \end{mat}
       \end{equation*}  
       为了找一个相伴的基，我们计算零空间。
       \begin{equation*}
         \nullspace{t+i}=\set{\colvec{-iy \\ y}
                                        \suchthat y\in\C} 
         \qquad
         \nullspace{t-i}=\set{\colvec{iy \\ y}
                                        \suchthat y\in\C} 
       \end{equation*}
       例如，
       \begin{equation*}
         T+i\cdot I=
         \begin{mat}[r]
           i  &-1  \\
           1  &i
         \end{mat}
       \end{equation*}
       于是通过求解这个线性方程组，
       我们得到 $t+i$ 的零空间的描述。
       \begin{equation*}
         \begin{linsys}{2}
           ix  &-  &y  &=  &0  \\
            x  &+  &iy &=  &0
         \end{linsys}
         \grstep{i\rho_1+\rho_2}
         \begin{linsys}{2}
           ix  &-  &y  &=  &0  \\
               &   &0  &=  &0
         \end{linsys}
       \end{equation*}
       （要把关系式 $ix=y$ 改写为用自由变量 $y$ 表示首变量 $x$
       的形式，可以将两
       边乘以 $-i$。）

       结果，其中一个这样的基如下。
       \begin{equation*}
         B=\sequence{\colvec[r]{-i \\ 1},
                     \colvec[r]{i \\ 1}}
       \end{equation*}  
     \end{answer}
  \item 
    对于仅有的特征值为 \( -3 \) 与 \( 4 \) 的
    \( \nbyn{3} \) 矩阵，有多少个相似类？
    \begin{answer}
     我们可以通过数可能的规范代表元，也就是数可能的
     Jordan 形式矩阵，来数出可能的类的个数。
     特征多项式必定是 $c_1(x)=(x+3)^2(x-4)$
     或者 $c_2(x)=(x+3)(x-4)^2$。
     在 $c_1$ 的情形，$t+3$ 在 $\gennullspace{t+3}$ 的一个链基上
     有两种可能的作用。
     \begin{equation*}
       \vec{\beta}_1\mapsto\vec{\beta}_2\mapsto \zero
       \qquad
       \begin{array}[t]{l}
         \vec{\beta}_1\mapsto\zero \\
         \vec{\beta}_2\mapsto\zero 
       \end{array}
     \end{equation*}
     有两个相伴的 Jordan 形式矩阵。
     \begin{equation*}
        \begin{mat}[r]
          -3  &0  &0  \\
           1  &-3 &0  \\
           0  &0  &4
        \end{mat}
        \qquad
        \begin{mat}[r]
          -3  &0  &0  \\
           0  &-3 &0  \\
           0  &0  &4
        \end{mat}
      \end{equation*}
      类似地，也有两个可能由 $c_2$ 产生的
      Jordan 形式矩阵。
      \begin{equation*}
        \begin{mat}[r]
          -3  &0  &0  \\
           0  &4  &0  \\
           0  &1  &4
        \end{mat}
        \qquad
        \begin{mat}[r]
          -3  &0  &0  \\
           0  &4  &0  \\
           0  &0  &4
        \end{mat}
     \end{equation*}  
     所以总共有四种可能的 Jordan 形式。
    \end{answer}
  \recommended \item
    证明：一个矩阵可对角化当且仅当它的极小
    多项式只有一次因式。
    \begin{answer}
       Jordan 形式是唯一的。
       对角矩阵本身就是 Jordan 形式。
       因此可对角化矩阵的 Jordan 形式就是它的对角化。
       若极小多项式有幂次高于一次的因式，
       则 Jordan 形式在次对角线上有 \( 1 \)，从而
       不是对角的。  
     \end{answer}
  \item 
     给出一个向量空间上的线性变换的例子，它没有非平凡的
     不变子空间。
     \begin{answer}
       一个例子是 \( \C \) 上把 \( x \) 映到 \( -x \) 的
        变换。  
      \end{answer}
  \item 
     证明：一个子空间是 \( t-\lambda_1 \) 不变的当且仅当
     它是 \( t-\lambda_2 \) 不变的。
     \begin{answer}
       两次应用 \nearbylemma{le:tInvIfftMinLambdaInv}；
       该子空间是 $t-\lambda_1$~不变的当且仅当它是
       $t$~不变的，而这又当且仅当它是
       $t-\lambda_2$~不变的。  
     \end{answer}
  \item 
     证明或否定：两个 \( \nbyn{n} \) 矩阵
     相似当且仅当它们有相同的特征多项式与
     极小多项式。
     \begin{answer}
       不成立；这两个 $\nbyn{4}$ 矩阵都有 $c(x)=(x-3)^4$
       与 $m(x)=(x-3)^2$。
       \begin{equation*}
          \begin{mat}[r]
             3  &0  &0  &0  \\
             1  &3  &0  &0  \\
             0  &0  &3  &0  \\
             0  &0  &1  &3
          \end{mat}
          \quad
          \begin{mat}[r]
             3  &0  &0  &0  \\
             1  &3  &0  &0  \\
             0  &0  &3  &0  \\
             0  &0  &0  &3
          \end{mat}
       \end{equation*} 
     \end{answer}
  \item 
    方阵的\definend{迹}\index{trace}\index{matrix!trace} 
    是其对角元之和。
    \begin{exparts}
       \partsitem 求出
         $\nbyn{2}$ 矩阵的特征多项式的公式。
       \partsitem 证明
         迹在相似下不变，从而我们可以有意义地
         谈论「映射的迹」。\index{linear map!trace}
         （\textit{提示：} 参见前面的题。）
       \partsitem 迹在矩阵等价下不变吗？
       \partsitem 证明映射的迹是其特征值之和
         （计入重数）。
       \partsitem 证明幂零映射的迹为零。
         逆命题成立吗？
    \end{exparts}
    \begin{answer}
      \begin{exparts}
         \partsitem 特征多项式如下。 
           \begin{equation*}
             \begin{vmat}
               a-x  &b  \\
               c  &d-x
             \end{vmat}
             =(a-x)(d-x)-bc=ad-(a+d)x+x^2-bc
             =x^2-(a+d)x+(ad-bc)
           \end{equation*}
           注意行列式作为常数项出现。
         \partsitem 回忆特征多项式
            \( \deter{T-xI} \) 在相似下不变。
            利用置换展开公式证明迹
            是 \( x^{n-1} \) 的系数的相反数。
         \partsitem 不，存在矩阵 $T$ 与 $S$，它们
            等价 $S=PTQ$（对某个非奇异的 $P$ 与 $Q$），
            但迹不同。
            一个简单的例子如下。
            \begin{equation*}
               PTQ=
               \begin{mat}[r]
                  2  &0  \\
                  0  &1
               \end{mat}
               \begin{mat}[r]
                  1  &0  \\
                  0  &1
               \end{mat}
               \begin{mat}[r]
                  1  &0  \\
                  0  &1
               \end{mat}
               =
               \begin{mat}[r]
                  2  &0  \\
                  0  &1
               \end{mat}
            \end{equation*}
            用 $\nbyn{1}$ 矩阵甚至能得到更简单的例子。
         \partsitem 把矩阵化成 Jordan 形式。
            由第一小题，迹不变。
         \partsitem 第一部分容易；用第三小题。
            逆命题不成立：~这个矩阵
            \begin{equation*}
               \begin{mat}[r]
                  1  &0  \\
                  0  &-1
               \end{mat}
            \end{equation*}
            迹为零但不是幂零的。
       \end{exparts}  
     \end{answer}
  \item 
    用 \nearbydefinition{def:invariant} 来检查一个子空间
    是否 $t$~不变时，我们似乎不得不检查（非平凡）子空间中无穷多个
    向量，看它们是否满足条件。
    证明：一个子空间是 \( t \)~不变的当且仅当其部分基具有如下性质：对其中所有元素，
    $t(\vec{\beta})$ 都在
    该子空间中。
    \begin{answer}
      设 \( B_M \) 是某个
      向量空间的子空间 \( M \) 的一个基。
      一个方向的蕴含是显然的；若 \( M \) 是 \( t \) 不变的，则
      特别地，若 \( \vec{m}\in B_M \)，那么 \( t(\vec{m})\in M \)。
      对于另一个方向，设
      \( B_M=\sequence{\vec{\beta}_1,\dots,\vec{\beta}_q} \)，注意
      \( t(\vec{m})=t(m_1\vec{\beta}_1+\dots+m_q\vec{\beta}_q)
                   =m_1t(\vec{\beta}_1)+\dots+m_qt(\vec{\beta}_q) \)
      在 \( M \) 中，因为任何子空间对线性
      组合封闭。 
    \end{answer}
  \recommended \item 
    \( t \) 不变性对交运算保持吗？
    对并呢？
    对补集呢？
    对子空间的和呢？
    \begin{answer}
      是的，$t$ 不变子空间的交
      是 $t$~不变的。
      设 \( M \) 与 \( N \) 都是 \( t \) 不变的。
      若 \( \vec{v}\in M\intersection N \)，则 \( t(\vec{v})\in M \) 
      由 \( M \) 的不变性有 \( t(\vec{v})\in M \)，由
      \( N \) 的不变性有 \( t(\vec{v})\in N \)。

      当然，两个子空间的并未必是子空间
      （记住 $x$ 轴\hbox{} 与 $y$ 轴是平面
      $\Re^2$ 的子空间，但两轴的并对向量加法不封闭；
      例如它不包含
      $\vec{e}_1+\vec{e}_2$。）
      然而，不变子集的并是不变子集；若
      \( \vec{v}\in M\union N \)，则 \( \vec{v}\in M \) 或 \( \vec{v}\in N \)，
      于是 \( t(\vec{v})\in M \) 或 \( t(\vec{v})\in N \)。

      不，不变子空间的补未必是不变的。
      考虑零变换之下 \( \C^2 \) 的子空间
      \begin{equation*}
        \set{\colvec{x \\ 0}\suchthat x\in\C}
      \end{equation*}
      。

      是的，两个不变子空间的和是不变的。
      验证很容易。  
    \end{answer}
  \item  \label{exer:OrderC}
     若某些特征值是复数，给出排列
     Jordan 块的一种方式。
     也就是说，为复数建议一种合理的顺序。
     \begin{answer}
       一种这样的顺序是\definend{字典序}。
       先按实部排序，再按 \( i \) 的系数排序。
       例如，\( 3+2i<4+1i \) 但 \( 4+1i<4+2i \)。  
     \end{answer}
  \item
    设 \( \polyspace_j(\Re) \) 是实数域上
    次数 \( j \) 的多项式构成的向量空间。
    证明：若 \( j\le k \)，则 \( \polyspace_j(\Re) \) 是
    \( \polyspace_k(\Re) \) 在微分算子下的不变子空间。
    在 \( \polyspace_7(\Re) \) 中，\( \polyspace_0(\Re) \)、
    \ldots、\( \polyspace_6(\Re) \) 中有哪一个有不变的补吗？
    \begin{answer}
     前一半容易\Dash 任何实系数多项式的导数是次数更低的
      实系数多项式。
      后一半的答案是「没有」；\( \polyspace_j(\Re) \) 的任何补
      必须包含一个次数为 \( j+1 \) 的多项式，
      而该多项式的导数在 \( \polyspace_j(\Re)\) 中。  
     \end{answer}
  \item
    在 \( \polyspace_n(\Re) \)（实数域上
    次数 \( n \) 的多项式构成的向量空间）中，
    \begin{equation*}
      \mathcal{E}=
      \set{p(x)\in\polyspace_n(\Re)\suchthat p(-x)=p(x) \text{\ 对所有\ }x}
    \end{equation*}
    与
    \begin{equation*}
      \mathcal{O}=
      \set{p(x)\in\polyspace_n(\Re)\suchthat p(-x)=-p(x) \text{\ 对所有\ }x}
    \end{equation*}
    分别是\definend{偶多项式}\index{even polynomials}\index{polynomial!even} 
    与\definend{奇多项式}\index{even polynomials}\index{polynomial!even}
     ；\( p(x)=x^2 \) 是
    偶的而 \( p(x)=x^3 \) 是奇的。
    证明它们是子空间。
    它们互补吗？
    它们在微分变换下不变吗？
    \begin{answer}
      对于前一半，证明每一个都是子空间，然后观察
      到任何多项式都可以唯一地
      写成偶次项与奇次项之和（零
      多项式两者都是）。
      后一半的答案是「不」：\( x^2 \) 是偶的，而
      \( 2x \) 是奇的。  
     \end{answer}
  % Lemma gone when moved to new JF presentation
  % \item 
  %   \nearbylemma{le:InvCompSubspSplitTrans} says that if \( M \) and
  %   \( N \) are
  %   invariant complements then \( t \) has a representation in the given
  %   block form (with respect to the same ending as starting basis, of course).
  %   Does the implication reverse?
  %   \begin{answer}
  %     Yes.
  %     If \( \rep{t}{B,B} \) has the given block form, take \( B_M \) to
  %     be the first \( j \) vectors of \( B \), where \( J \) is the
  %     \( \nbyn{j} \) upper left submatrix.
  %     Take \( B_N \) to be the remaining \( k \) vectors in \( B \).
  %     Let \( M \) and \( N \) be the spans of \( B_M \) and \( B_N \).
  %     Clearly \( M \) and \( N \) are complementary.
  %     To see \( M \) is invariant (\( N \) works the same way), represent
  %     any \( \vec{m}\in M \) with respect to \( B \), note the last
  %     \( k \) components are zeroes, and multiply by the given block
  %     matrix.
  %     The final \( k \) components of the result are zeroes, so that
  %     result is again in \( M \). 
  %    \end{answer}
   \item 
     若矩阵 \( S \) 满足 \( S^2=T \)，
     则称 \( S \) 是 \( T \) 的\definend{平方根}\index{square root}。
     证明任何非奇异矩阵都有平方根。
     \begin{answer}
         把矩阵化成 Jordan 形式。
         由非奇异性，对角线上没有零特征值。
         仿照这个例子：
          \begin{equation*}
             \begin{mat}[r]
                9  &0  &0 \\
                1  &9  &0 \\
                0  &0  &4
             \end{mat}
             =
             \begin{mat}[r]
                3  &0  &0 \\
               1/6 &3  &0 \\
                0  &0  &2
             \end{mat}^2
          \end{equation*}
         来构造一个平方根。
         证明它在相似下仍然成立：若 \( S^2=T \)，则
         \( (PSP^{-1})(PSP^{-1})=PTP^{-1} \)。 
    \end{answer}
\index{Jordan form|)}
\end{exercises}
